\section{Introduction}
Gas-phase kinetic studies underpin our understanding of key atmospheric processes such as stratospheric ozone formation, regulation and depletion \cite{Rowland2006, Dameris2010}, and air pollution (e.g., \ce{NO_x}, tropospheric ozone, peroxy acetyl nitrate (PAN), and aerosol) generation and reduction \cite{VonSchneidemesser2015, Stockwell2020}. Research into this field is thus crucial for air quality and climate change assessment and mitigation. The aircraft sector is an important component in the assessment of air pollution and climate change, as it is one of only two non-surface sources of \ce{NO_x} emissions \cite{Wasiuk2016, Tait2022, Schumann1997}, with the other being lightning \cite{Tie2002, Schumann2007, Stockwell1999}. In addition to the direct emission of \ce{CO_2} through jet fuel combustion, aircraft emit \ce{NO_x} which contributes to tropospheric ozone formation, as well as water vapour and particulates that give rise to contrail formation \cite{Lee2021}. 

Traditionally, mitigation efforts have gone toward incremental increases in efficiency, due to technological advancements and innovative design. However, due to the increasing urgency to rapidly reduce aviation climate impact, the focus has shifted toward more radical measures. This includes the development of new propulsion systems (e.g., liquid hydrogen) \cite{Yusaf2022, Khan2022}, sustainable aviation fuels (SAFs) \cite{Doliente2020}, and changes to aircraft flight routes and operations. Unlike the former two methods, modifications to flight routing pose the most viable solution for implementation this decade \cite{Niklass2019a}. This paper explores the potential of formation flying to reduce the ozone production efficiency of aircraft \ce{NO_x} emissions. 

Reducing \ce{NO_x} emissions requires a lower temperature burn and this conflicts with trends toward higher temperature, more energy efficient jet fuel combustion. Formation flight may hold the key to maximizing the reduction in ozone production, due to any additional \ce{NO_x} that is output from aircraft. The concept involves two or more aircraft flying no more than 1--2 km apart longitudinally, akin to energy-saving migratory flocks of geese. This enables follower aircraft to “surf” the wake of their leader, thus reducing required lift and hence thrust and fuel burn, on the order of 5\% per aircraft \cite{Xu2014, Bower2009}.

Further to the fuel saving achieved through this concept, there are also two notable chemical saturation effects that occur, due to the local accumulation of aircraft emissions. First, in airspace conditions conducive to persistent contrail formation, formation flight can help reduce the net contrail effect. This is because the generation of two or more persistent contrails in the same region will compete for available water vapour. Thus, they mutually inhibit growth and the combined warming effect is lessened compared with if they were generated individually. Unterstrasser et al. \cite{Unterstrasser2020} explores this effect in great detail. 

The second saturation effect to be observed is the accumulation of \ce{NO_x} emissions, leading to \ce{NO_x}-saturated conditions and a negative correlation with ozone production efficiency. It is well known in the urban environment, where there can be very high \ce{NO_x} levels, that ozone formation is virtually zero as odd hydrogen (\ce{HO_x}) radical concentrations become significantly depleted. In this environment, \ce{NO_2} favours the formation of nitric acid (\ce{HNO3}) over ozone \cite{Jenkin2002, Sillman2000}. \ce{HNO3} is eventually incorporated into aerosol and removed from the atmosphere, representing a permanent sink for \ce{NO_x} and preventing ozone formation. Can formation flying also deliver a reduction in ozone formation following \ce{NO_x} emission? If \ce{NO_x} levels are maintained at high enough levels, OH collapse could occur, resulting in a significant fraction of NO2 being sequestered to form \ce{HNO3} and therefore being unable to contribute to ozone formation. Formation flight scenarios can provide climate benefits, due to this \ce{NO2} loss mechanism reducing net ozone formation (compared to the same emission plumes released independently of one another). In essence, when emitted, \ce{NO_x} can, under high enough VOC levels, catalyse the formation of ozone. This can be demonstrated through the CO oxidation cycle, using it as a representative VOC, even though it is not strictly speaking one:

\reaction[OH+CO_4]{OH + CO -> H + CO_2}

\reaction[H+O2+M_4]{H + O_2 + M -> HO_2 + M}

\reaction[HO2+NO_4]{HO_2 + NO -> OH + NO_2}

\reaction[NO2+hv_4]{NO_2 + hv -> NO + O}

\reaction[O+O2+M_4]{O + O_2 + M -> O_3 + M}

Through reactions \eqref{OH+CO_4}--\eqref{O+O2+M_4}, ozone is produced through \ce{NO_2} photolysis, and the longer the chain length converting NO to \ce{NO_2} \eqref{HO2+NO_4} and back to NO \eqref{NO2+hv_4} the more ozone is formed. However, if NO and \ce{NO_2} levels are high enough, the reaction of \ce{NO_2} with OH (reaction \eqref{OH+NO2+M_4}) can start to compete with the reaction of OH with CO and \ce{CH_4} in the mid to upper troposphere (reactions \eqref{OH+CO_4} and \eqref{OH+CH4_4}).

\reaction[OH+NO2+M_4]{OH + NO_2 + M -> HNO3 + M}

\reaction[OH+CH4_4]{OH + CH4 -> products}

The reason why reaction \eqref{OH+NO2+M_4} may start to compete in the upper troposphere, where temperatures are far lower than they are at the surface, is that as an association reaction, the rate coefficient of the reaction \eqref{OH+NO2+M_4}, $k_6$, increases with decreasing temperature \cite{Atkinson2004, Butkovskaya2005} whereas rate coefficient of the reaction \eqref{OH+CH4_4}, $k_7$ decreases with decreasing temperature \cite{Atkinson2003} in a typical Arrhenius type behavior. Rate constant of the reaction \eqref{OH+CO_4}, $k_1$ has both a direct and association channel but appears to vary very little with temperature.

Aircraft \ce{NO_x} perturbs the fast photochemistry of the UTLS by interfering with the balance between photochemical ozone production and ozone destruction. This balance is controlled by the \ce{NO_x}-compensation point, the point at which with decreasing \ce{NO_x} levels, photochemical ozone production through reaction \eqref{HO2+NO_4} exactly balances photochemical ozone destruction through:

\reaction[HO2+O3_4]{HO_2 + O_3 -> OH + O_2 + O_2}

In the UTLS region, the \ce{NO_x} compensation point, which is set by the rate coefficients for the reactions \eqref{HO2+NO_4} and \eqref{HO2+O3_4}, is about 10~ppt as NO (calculated assuming the rate coefficient of reaction \eqref{HO2+NO_4}, $k_3$ of $1.19 \times 10^{-11} \text{cm}^3 \text{molecule}^{-1} \text{s}^{-1}$, the rate coefficient of reaction \eqref{HO2+O3_4}, $k_8$ of $1.13 \times 10^{-15} \text{cm}^3 \text{molecule}^{-1} \text{s}^{-1}$ based on Atkinson et al. \cite{Atkinson2004} and \ce{O_3} 100~ppb). That is to say, if the NO mixing ratio is less than 10~ppt, the fast \ce{HO_x} photochemistry leads to ozone destruction. This is the general state of the UTLS region under typical conditions of temperature and pressure. In the presence of aircraft at cruise altitude, through their \ce{NO_x} emissions, the \ce{NO_x} compensation point may be exceeded and photochemical ozone production sets in, leading ultimately to an increase in global warming since the cold conditions of the UTLS favour enhanced radiative forcing from the increased ozone mixing ratios. The elevated OH concentrations driven by reaction \eqref{HO2+NO_4} lead to elevated \ce{CH_4} destruction and hence decreased build-up, with its attendant negative radiative forcing and global cooling.

Most of the studies of the impacts of aircraft cruise emissions have been completed using sophisticated three-dimensional (3D) Eulerian grid-based chemistry-transport models (CTMs). Despite a high degree of sophistication, typical spatial resolutions are no better than 1~\textdegree $\times$ 1~\textdegree (say, 100~km $\times$ 100~km in the North Atlantic air traffic corridor). A review of these studies is available elsewhere \cite{Fuglestvedt1999}. Current CTM studies of these greenhouse gas impacts may well overstate the impacts of aircraft cruise emissions. This arises because of the current mismatch between the coarse spatial resolution of the CTMs and the fine-scale dimensions of the major air traffic corridors. A similar conclusion may be reached by consideration of the temporal resolution of these studies. The CTM studies may necessarily average aircraft cruise emissions over months or even years whereas on average flights last hours \cite{Wasiuk2016}. Together, the lack of spatial and temporal CTM resolution may ensure that aviation \ce{NO_x} is present in excess of the \ce{NO_x} compensation point at all times and in all grid squares through which a significant number of aircraft fly.

In this paper, first a kinetic model feasibility study was carried out to determine what \ce{NO_x} threshold level is required to start reducing ozone formation. Second, an established 3D global CTM, STOCHEM-CRI, was used to identify regions of the atmosphere where the lowest \ce{NO_x} thresholds exist. Thirdly, validation of the results from the global analysis of \ce{NO_x} threshold concentrations was conducted by running the simulations with a much higher resolution dataset, and comparing these results with that derived from coarse STOCHEM-CRI model. The In-service Aircraft for Global Observing System (IAGOS) data \cite{IAGOS} is used in place of the 5~\textdegree $\times$ 5~\textdegree monthly global model fields from STOCHEM-CRI \cite{Khan2021}. Finally, we develop a simple aircraft plume dispersion model which addresses the issues of inadequate spatial and temporal CTM resolution and understands how different the chemistry is likely to look under conditions of higher spatial and temporal resolution.